\documentclass[11pt,a4paper]{article}
\usepackage[T1]{fontenc}
\usepackage{isabelle,isabellesym}
\usepackage{amssymb}
\usepackage{pdfsetup}

% urls in roman style, theory text in math-similar italics
\urlstyle{rm}
\isabellestyle{it}


\begin{document}

\title{The Five Platonic Solids}
\author{Evan Finken}
\maketitle

\begin{abstract}
  This entry formalizes that there are exactly five Platonic solids: the
  tetrahedron, cube, octahedron, dodecahedron, and icosahedron. A Platonic
  solid is a convex polyhedron whose faces are congruent regular polygons,
  with the same number of edges meeting at each vertex. Euler's polyhedron
  formula is used to show that there are at most five Platonic solids, and
  each of the five, defined as the convex hull of its vertices, is then
  shown to satisfy the required properties. The formal proof is computational (slow!)\ 
  and is ported from John Harrison's formalization in HOL Light~\cite{holLightPlatonic}.
  Large language models were used to explain the intermediate steps of the
  HOL Light formalization and to write parts of the more complicated proofs.
\end{abstract}

\tableofcontents
\newpage

\input{session}

\bibliographystyle{abbrv}
\bibliography{root}

\end{document}
